在过去的学习中，我们已经学习了许多物理定律及其应用。对于一大类问题，通过已知定律
并结合系统分析，已经可以得到结果，其中也包括一些相当复杂的问题。但在学习过程中会
发现，处理某些问题时会遇到数学上的困难。例如，位置随时间的变化率被定义为速度，可
我们还没有掌握由任意给定运动求速度的普遍方法；又如，速度—时间图像与横轴所围面积
表示位移，却还不能对一般的速度函数求任意一段时间内的位移。

历史上，人们在研究引力及其性质时也遇到了类似困难。牛顿（I. Newton）和莱布尼茨
（G. W. Leibniz）分别独立地发明了新的数学工具——微积分（calculus），使这类问题的
处理成为可能。微积分及随之发展起来的数学理论极大地推动了物理学的发展，不仅使我们
能够处理许多与曲线、曲面性质有关的问题，其中也包含着重要的物理思想。

具有光滑图像的单自变量函数在某点处切线的斜率称为函数在该点的\textbf{导数}
（derivative）。若函数记为 $y=y(x)$，曲线在 $x$ 点处的切线可以看作割线的极限，其
斜率为
\begin{equation}
 y'(x)=\lim_{\Delta x\to0}\frac{y(x+\Delta x)-y(x)}{\Delta x}.
\end{equation}
导数有多种表示方法，例如 $y'(x)$、$\dfrac{\dd y}{\dd x}(x)$。物理中遇到的绝大部分
函数都可以求导；除了会计算特定函数的导数，还需要理解导数的物理意义。

直线运动质点的速度是坐标对时间的导数：
\begin{equation}
 v(t)=\frac{\dd x}{\dd t}(t),
\end{equation}
加速度则是速度对时间的导数：
\begin{equation}
 a(t)=\frac{\dd v}{\dd t}(t)=\frac{\dd^2x}{\dd t^2}(t).
\end{equation}
因此加速度也是坐标对时间的二阶导数（second order derivative），常记作 $x''$ 或
$\dfrac{\dd^2x}{\dd t^2}$；类似地还可以定义更高阶导数。

许多物理量本来就是借助导数定义的。例如实验给出一般弹性体的外力与形变量的函数
$F(x)$，它在形变 $x$ 处的切线刚度（局部劲度）可写成
\begin{equation}
 k(x)=\frac{\dd F}{\dd x}(x).
\end{equation}
胡克定律只是切线刚度为常数的特殊情况。又如实验测得某物体吸收的热量与温度之间的关系
$Q(T)$，则该物体在温度 $T$ 处的热容量为
\begin{equation}
 C(T)=\frac{\dd Q}{\dd T}(T).
\end{equation}
对于真实物体，热容量一般会随温度改变。

利用导数的定义和基本运算法则，可以求出物理中绝大部分常见函数的导数。下面通过例题
熟悉基本求导规律。

\subsection{局部变化率、链式法则与偏导数}

前述割线斜率的极限说明，导数不是“两个无穷小量的普通商”，而是一个局部变化率。位置对时间的一阶导数是速度，二阶导数是
加速度；电量对时间的导数是电流；弹性体的切线刚度可写成 $\dd F/\dd x$；若物体吸热量为
$Q(T)$，则温度为 $T$ 时的热容 $C(T)$ 为 $\dd Q/\dd T$。这些定义都在问同一件事：系统处于某个状态时，
再作极小改变会得到怎样的即时响应。

复合关系要用链式法则。若 $y=f(u)$、$u=g(x)$，则
\[
 \frac{\dd y}{\dd x}=\frac{\dd y}{\dd u}\frac{\dd u}{\dd x}.
\]
例如势能 $U(x)$ 随位置变化，而位置又随时间变化，就有
$\dd U/\dd t=(\dd U/\dd x)v$。若函数由多个物理量共同决定，则应说明求偏导时保持哪些量不变；
热学中 $(\partial p/\partial V)_T$ 与 $(\partial p/\partial V)_S$ 就代表不同的物理过程。

\subsection{极值判别、物理约束与近似展开}

若函数在定义域内部的 $x_0$ 处可导并取得局部极值，则必有 $f'(x_0)=0$；但驻点不一定是极值，
还要比较左右导数或二阶导数。若 $f'(x_0)=0$ 且 $f''(x_0)>0$，则 $x_0$ 是严格局部极小点；
若 $f'(x_0)=0$ 且 $f''(x_0)<0$，则 $x_0$ 是严格局部极大点。端点和不可导点也可能
给出最大、最小值，故完整求极值必须连同定义域一起检查。物理答案还要满足长度为正、接触不脱离、
速度可达等约束，纯代数驻点未必可实现。

在 $x_0$ 附近，足够光滑的函数可逐阶展开：
\[
 f(x_0+\xi)=f(x_0)+f'(x_0)\xi+
 \frac12f''(x_0)\xi^2+\frac1{3!}f'''(x_0)\xi^3+\cdots .
\]
近似保留到哪一阶取决于所需精度。若 $x_0$ 是势能的平衡位置，一阶项消失，最低非零项通常是二次项，
系统因而近似简谐振动；若二次项也为零，就必须保留更高阶，回复力和振动周期会表现出不同规律。
小角近似 $\sin\theta\simeq\theta$、$\cos\theta\simeq1-\theta^2/2$ 也来自同一展开，使用前要说明
$|\theta|\ll1$ 且角度以弧度计。

对时间求导时还常使用点记号：$\dot X=\dd X/\dd t$，
$\ddot X=\dd^2X/\dd t^2$；对一般自变量仍写撇号或 $\dd/\dd x$。
导数还提供变化率的量级判断。设某量的特征值为 $X$、变化时间为 $T$，则
$\dot X$ 的量级约为 $X/T$；两个机制谁占主导，可比较相应导数项的无量纲比值。竞赛推导中，
这类先估量级、再保留主导项的做法，常比直接处理完整方程更清楚。


\begin{origexamplebox}{例 1.1}
根据定义求下列函数的导数：
\[
 y(x)=C,\qquad y(x)=kx,\qquad y(x)=ax^n,\qquad
 y(x)=A\sin(kx),\qquad y(x)=\e^x,
\]
其中除自变量 $x$ 外，其余字母均为常数；幂函数一式先取非负整数 $n$。
\end{origexamplebox}

\begin{xesolution}
直接使用导数定义。例如幂函数有
\[
 \frac{\dd}{\dd x}(ax^n)
 =a\lim_{\Delta x\to0}
 \frac{(x+\Delta x)^n-x^n}{\Delta x}.
\]
由有限二项式展开
\[
 (x+\Delta x)^n-x^n
 =nx^{n-1}\Delta x+\binom n2x^{n-2}(\Delta x)^2+\cdots+(\Delta x)^n,
\]
除以 $\Delta x$ 后，第一项为 $nx^{n-1}$，其余各项都至少含一个 $\Delta x$，在
$\Delta x\to0$ 时趋于零，故得到 $anx^{n-1}$。对实数指数，在 $x>0$ 的区间可由广义二项式展开得到同一公式。同理可得
\[
 \boxed{C'=0},\qquad
 \boxed{(kx)'=k},\qquad
 \boxed{(ax^n)'=anx^{n-1}},
\]
\[
 \boxed{[A\sin(kx)]'=Ak\cos(kx)},\qquad
 \boxed{(\e^x)'=\e^x}.
\]

下面把三角函数和指数函数两式的极限过程完整写出。对 $y=A\sin(kx)$，由导数定义
\begin{align*}
 \frac{\dd y}{\dd x}
 &=A\lim_{\Delta x\to0}
 \frac{\sin[k(x+\Delta x)]-\sin(kx)}{\Delta x}\\
 &=A\lim_{\Delta x\to0}
 \frac{2\cos\!\left(kx+\dfrac{k\Delta x}{2}\right)
       \sin\!\left(\dfrac{k\Delta x}{2}\right)}{\Delta x}.
\end{align*}
第二步使用了恒等式
$\sin\alpha-\sin\beta
=2\cos[(\alpha+\beta)/2]\sin[(\alpha-\beta)/2]$。
把分式改写成
\[
 Ak\cos\!\left(kx+\frac{k\Delta x}{2}\right)
 \frac{\sin(k\Delta x/2)}{k\Delta x/2}.
\]
当 $\Delta x\to0$ 时，余弦因子趋于 $\cos(kx)$，而令
$u=k\Delta x/2$ 后，需要使用基本极限 $\sin u/u\to1$。这个极限可由单位圆直接证明：
当 $0<u<\pi/2$ 时，扇形面积与内、外两个三角形面积的比较给出
\[
 \sin u<u<\tan u.
\]
各项除以 $u$ 并整理，得到
\[
 \cos u<\frac{\sin u}{u}<1.
\]
当 $u\to0^+$ 时两侧都趋于 $1$，由夹逼定理可得右极限为 $1$；又因为
$\sin(-u)/(-u)=\sin u/u$，左极限相同。因此
\[
 \frac{\dd}{\dd x}[A\sin(kx)]=Ak\cos(kx).
\]

对 $y=\e^x$，先把公因子 $\e^x$ 提出：
\begin{align*}
 \frac{\dd}{\dd x}\e^x
 &=\lim_{\Delta x\to0}
 \frac{\e^{x+\Delta x}-\e^x}{\Delta x}\\
 &=\e^x\lim_{\Delta x\to0}
 \frac{\e^{\Delta x}-1}{\Delta x}.
\end{align*}
把自然指数函数写成它的收敛级数定义
\[
 \e^u=1+u+\frac{u^2}{2!}+\frac{u^3}{3!}+\cdots,
\]
便有
\[
 \frac{\e^u-1}{u}
 =1+\frac{u}{2!}+\frac{u^2}{3!}+\cdots.
\]
当 $|u|<1/2$ 时，余项的绝对值满足
\[
 \left|\frac{u}{2!}+\frac{u^2}{3!}+\cdots\right|
 \le |u|+|u|^2+\cdots=\frac{|u|}{1-|u|}\longrightarrow0.
\]
因此 $\lim_{u\to0}(\e^u-1)/u=1$。取 $u=\Delta x$，上式中的极限等于 $1$，最终得到
$\dd(\e^x)/\dd x=\e^x$。
\end{xesolution}

\begin{origexamplebox}{例 1.2}
函数 $f(x),g(x)$ 对 $x$ 的导数已知。根据导数定义证明以下求导运算的基本规律：
\[
 \frac{\dd}{\dd x}[a f(x)]=a\frac{\dd f}{\dd x},
\]
\[
 \frac{\dd}{\dd x}[f(x)g(x)]
 =f\frac{\dd g}{\dd x}+g\frac{\dd f}{\dd x},
\]
\[
 \frac{\dd}{\dd x}f(g(x))=\frac{\dd f}{\dd g}\frac{\dd g}{\dd x}.
\]
\end{origexamplebox}

\begin{xesolution}
对第一式，常数 $a$ 可直接移到极限号外：
\[
 \lim_{\Delta x\to0}\frac{af(x+\Delta x)-af(x)}{\Delta x}
 =a f'(x).
\]
对乘积，在分子中加上再减去 $f(x+\Delta x)g(x)$：
\begin{align*}
 &\frac{f(x+\Delta x)g(x+\Delta x)-f(x)g(x)}{\Delta x}\\
 &\quad=f(x+\Delta x)\frac{g(x+\Delta x)-g(x)}{\Delta x}
 +g(x)\frac{f(x+\Delta x)-f(x)}{\Delta x}.
\end{align*}
令 $\Delta x\to0$，由连续性 $f(x+\Delta x)\to f(x)$，便得乘积法则。

对复合函数令 $\Delta g=g(x+\Delta x)-g(x)$。当 $\Delta g\ne0$ 时
\[
 \frac{f(g+\Delta g)-f(g)}{\Delta x}
 =\frac{f(g+\Delta g)-f(g)}{\Delta g}\,
  \frac{\Delta g}{\Delta x}.
\]
不过 $\Delta g$ 可能在某些 $\Delta x\ne0$ 时恰为零，不能把上式中的除法当成总能进行。
为同时包含这个分支，利用 $f$ 在 $g(x)$ 处可导，把它的增量写为
\[
 f(g+\Delta g)-f(g)=f'(g)\Delta g+\varepsilon(\Delta g)\Delta g,
 \qquad \varepsilon(h)\to0\quad(h\to0).
\]
余项表示“实际增量减去切线估计”的差比 $\Delta g$ 更小；当 $\Delta g=0$ 时，
取 $\varepsilon(0)=0$，这条等式仍成立。只除以 $\Delta x\ne0$，得到
\[
 \frac{f(g+\Delta g)-f(g)}{\Delta x}
 =[f'(g)+\varepsilon(\Delta g)]\frac{\Delta g}{\Delta x}.
\]
$g$ 可导因而连续，所以 $\Delta x\to0$ 时 $\Delta g\to0$，
右侧两因子分别趋于 $f'(g(x))$ 和 $g'(x)$。这就严格得到链式法则
\[
 \boxed{\frac{\dd f(g(x))}{\dd x}=f'(g(x))g'(x)}.
\]
\end{xesolution}

\subsection{物理量的变化率}

导数在物理中最直接的应用，是求一个物理量对另一个物理量的变化率。这种计算广泛存在
于物理学的各个分支：速度是坐标随时间的变化率，加速度是速度随时间的变化率，保守力
也与势能随坐标的变化率直接相关。若能写出两个物理量之间的函数关系，直接求导便可
得到一个物理量改变时另一个物理量的变化率。

\begin{origexamplebox}{例 1.3}
一个一维运动质点的坐标随时间的关系分别为下列函数，求瞬时速度、加速度与时间的函数：
\begin{enumerate}[label=\arabic*.,leftmargin=2.2em]
 \item $x(t)=v_0t+\dfrac12at^2$；
 \item $x(t)=A\e^{-\beta t}\cos(\omega t)$，$\beta,\omega$ 为已知数；
 \item $x(t)=\dfrac{mv_0}{\beta}\left(1-\e^{-\frac{\beta}{m}t}\right)$，
       $\beta,m$ 为已知数。
\end{enumerate}
\end{origexamplebox}

\begin{xesolution}
逐式对 $t$ 求一阶、二阶导数。

\textbf{(1)}
\[
 \boxed{v(t)=v_0+at},\qquad \boxed{a(t)=a}.
\]

\textbf{(2)} 先使用乘积法则，再使用链式法则：
\[
 \boxed{v(t)=-A\e^{-\beta t}
   [\beta\cos(\omega t)+\omega\sin(\omega t)]},
\]
\[
 \boxed{a(t)=A\e^{-\beta t}
   [(\beta^2-\omega^2)\cos(\omega t)
    +2\beta\omega\sin(\omega t)]}.
\]

\textbf{(3)} 指数项中的 $-\beta/m$ 还要再乘一次：
\[
 \boxed{v(t)=v_0\e^{-\frac{\beta}{m}t}},\qquad
 \boxed{a(t)=-\frac{\beta v_0}{m}\e^{-\frac{\beta}{m}t}}.
\]
由前两式可得
\[
 a=\frac{\dd v}{\dd t}=-\frac{\beta}{m}v.
\]
在本题给出的运动中 $v>0$ 而 $a<0$，所以加速度方向与速度方向相反，速率随时间减小。
这种“与速度反向、使速率减小的加速度分量”常被称为减速度；它不是一种新的物理量，
仍然是加速度。若把这一运动解释为物体只受线性阻力，或其他恒力已经彼此平衡，则由
牛顿第二定律
\[
 F_{\rm net}=ma=-\beta v
\]
可知阻力为 $F_{\rm D}=-\beta v$：负号表示阻力与速度反向，大小 $\beta v$ 随速率增大。
若没有这一受力模型，仅由给定的 $x(t)$ 只能判断物体在减速，不能凭空断定具体是哪一种力造成的。
\end{xesolution}

\Needspace{0.38\textheight}
\begin{origexamplebox}{例 1.4}
一个半径为 $R$ 的半圆柱体在垂直于柱轴的水平方向上以匀速 $v_0$ 运动。如图，在圆柱面
上搁一竖直杆，此杆只能沿竖直方向上下运动。求杆上任意一点的速度、加速度与图示角
$\theta$ 的关系。
\probdiagOneFour
\end{origexamplebox}

\begin{xesolution}
设半圆柱圆心的横坐标为 $X$，杆的横坐标为 $x_{\rm r}$，接触点相对圆心的高度为 $y$。
由解图可见
\[
 x_{\rm r}=X+R\sin\theta,\qquad y=R\cos\theta.
\]
杆只能竖直运动，所以 $x_{\rm r}$ 为常数，而 $\dot X=v_0$。对第一式求导：
\[
 0=v_0+R\cos\theta\,\dot\theta,\qquad
 \dot\theta=-\frac{v_0}{R\cos\theta}.
\]
再对 $y$ 求导，得到杆的竖直速度
\[
 v_y=\dot y=-R\sin\theta\,\dot\theta
 =\boxed{v_0\tan\theta}.
\]
\soldiagOneFour
继续求导：
\[
 a_y=v_0\sec^2\theta\,\dot\theta
 =\boxed{-\frac{v_0^2}{R\cos^3\theta}}.
\]
若取竖直向上为正，速度向上，加速度向下。刚性杆作平动，所以杆上任意一点的速度和
加速度都相同。当 $\theta\to\pi/2$ 时理想模型给出发散结果，说明接触状态不可能一直
保持到圆柱侧缘。
\end{xesolution}

\begin{origexamplebox}{例 1.5}
一个发光点在焦距为 $f$ 的凸透镜主光轴上，以恒定不变的速度 $v$ 向着透镜运动。求其
像点移动速度与物距 $u$ 的关系。
\end{origexamplebox}

\begin{xesolution}
设物体在透镜左侧，$u>0$ 是它到透镜的距离；取像坐标 $q$ 向右为正，
所以右侧实像 $q>0$、左侧虚像 $q<0$。在近轴近似下，对焦距 $f>0$ 的薄凸透镜，
\[
 \frac1f=\frac1u+\frac1q,\qquad q=\frac{fu}{u-f}.
\]
光点向透镜运动，故 $\dot u=-v$。对 $q(u)$ 求导：
\[
 \frac{\dd q}{\dd u}=-\frac{f^2}{(u-f)^2},
\]
从而
\[
 \boxed{\dot q=\frac{f^2}{(u-f)^2}v}.
\]
当 $u>f$ 时，像位于透镜另一侧，$\dot q>0$ 表示实像远离透镜；物体越接近焦点，像点
移动越快。当 $u=f$ 时，出射光平行，不存在有限位置的像点，因此此时不能给像点一个有限速度。
当 $0<u<f$ 时，$q<0$，形成同侧虚像；虽然仍有 $\dot q>0$，
其含义却是负的像坐标向零增加，虚像向透镜靠近，而不是远离。
例如 $u=f/2$ 时 $q=-f$、$\dot q=4v$，这与有向坐标的判断一致。
不能把 $u=f$ 两侧的像当作以有限速度连续越过无穷远的同一个物点。
\end{xesolution}

\subsection{求极值}

对于光滑函数，函数曲线上各点的切线斜率就是该点的导数。当某一点处切线水平时，该点
可能是局部极大值、局部极小值，也可能既非极大又非极小。导数为零的点统称为\textbf{驻点}
（stationary point）。若在驻点处二阶导数为正，则是局部极小值；二阶导数为负，则是
局部极大值；二阶导数为零时，还需进一步判断。

\extremaDiagram

当感兴趣的物理量随其他因素变化时，先写出它们之间的函数关系，再令一阶导数为零，便
可寻找变化过程中的极值；必要时再借助二阶导数或端点条件判断极值类型。

\Needspace{0.38\textheight}
\begin{origexamplebox}{例 1.6}
利用费马原理证明光的折射定律：两个均匀介质的折射率分别为 $n_1,n_2$，分界面为一
平面。证明由平面外一点 $A$ 到达另一点 $B$ 的光线中，满足折射定律的光线的光程为
极值。
\probdiagOneSix
\end{origexamplebox}

\begin{xesolution}
设入射点 $P$ 到 $A$ 在界面上的垂足的距离为 $x$，两垂足间的水平距离为 $L$，而 $A,B$
到界面的垂直距离分别为 $h_1,h_2$。光程为
\[
 \mathcal L(x)=n_1\sqrt{h_1^2+x^2}
 +n_2\sqrt{h_2^2+(L-x)^2}.
\]
费马原理要求实际光路使光程取稳定值，因此
\[
 \frac{\dd\mathcal L}{\dd x}
 =n_1\frac{x}{\sqrt{h_1^2+x^2}}
 -n_2\frac{L-x}{\sqrt{h_2^2+(L-x)^2}}=0.
\]
由图中几何关系
\[
 \sin\alpha=\frac{x}{\sqrt{h_1^2+x^2}},\qquad
 \sin\beta=\frac{L-x}{\sqrt{h_2^2+(L-x)^2}},
\]
所以
\[
 \boxed{n_1\sin\alpha=n_2\sin\beta},
\]
这就是斯涅耳折射定律。此外
\[
 \frac{\dd^2\mathcal L}{\dd x^2}
 =\frac{n_1h_1^2}{(h_1^2+x^2)^{3/2}}
 +\frac{n_2h_2^2}{[h_2^2+(L-x)^2]^{3/2}}>0,
\]
故该稳定光路给出的确是光程的极小值。
\end{xesolution}

\Needspace{0.40\textheight}
\begin{origexamplebox}{例 1.7}
求理想气体在下面两个热力学过程中的温度最大值和最小值：
\begin{enumerate}[label=(\alph*),leftmargin=2.2em]
 \item 气体从状态 $(2p_0,V_0)$ 沿连接两点的直线段变化到状态 $(p_0,2V_0)$；
 \item $p$--$V$ 图中的一个圆（椭圆），过程满足
 \[
  \left(\frac p{p_0}-1\right)^2+
  \left(\frac V{V_0}-1\right)^2=\left(\frac12\right)^2.
 \]
\end{enumerate}
\probdiagOneSeven
\end{origexamplebox}

\begin{xesolution}
对一定量理想气体，$T=pV/(nR_g)$，因此只需在给定曲线上求乘积 $pV$ 的极值，其中
$R_g$ 为气体常量。

\textbf{(a)} 连接两状态的直线为
\[
 p=3p_0-\frac{p_0}{V_0}V.
\]
于是
\[
 pV=3p_0V-\frac{p_0}{V_0}V^2,\qquad
 \frac{\dd(pV)}{\dd V}=3p_0-2\frac{p_0}{V_0}V.
\]
在题设线段上 $V_0\le V\le2V_0$。驻点位于 $V=3V_0/2$，此时 $p=3p_0/2$，为抛物线的最大值；
最小值只能出现在两个端点，而两端点的温度相同，故
\[
 \boxed{T_{\max}=\frac{9p_0V_0}{4nR_g}},\qquad
 \boxed{T_{\min}=\frac{2p_0V_0}{nR_g}}.
\]

\textbf{(b)} 令 $P=p/p_0, X=V/V_0$，并参数化圆：
\[
 P=1+\frac12\cos\phi,\qquad X=1+\frac12\sin\phi.
\]
于是 $T=(p_0V_0/nR_g)PX$。求导可得
\[
 \frac{\dd(PX)}{\dd\phi}
 =\frac12(\cos\phi-\sin\phi)+\frac14\cos2\phi
 =(\cos\phi-\sin\phi)
 \left(\frac12+\frac14(\cos\phi+\sin\phi)\right).
\]
第二因子恒正，故驻点满足 $\sin\phi=\cos\phi$。比较两组对径点：
\[
 P=X=1\pm\frac{1}{2\sqrt2}.
\]
因此
\[
 \boxed{T_{\max}=\frac{p_0V_0}{nR_g}
  \left(1+\frac1{2\sqrt2}\right)^2},
\]
\[
 \boxed{T_{\min}=\frac{p_0V_0}{nR_g}
  \left(1-\frac1{2\sqrt2}\right)^2}.
\]
\end{xesolution}

\Needspace{0.30\textheight}
\begin{origexamplebox}{例 1.8}
宇宙中有两个天体，质量均为 $M$，相距 $2L$。求在它们连线的垂直平分线上何处的小
天体受到两个大天体的万有引力合力能够达到最大值。
\probdiagOneEight
\end{origexamplebox}

\begin{xesolution}
设小天体质量为 $m$，到连线中点的距离为 $x$。到每个大天体的距离为
$r=\sqrt{x^2+L^2}$。两个引力的垂直分量相消，沿垂直平分线方向的分量相加：
\[
 F(x)=2\frac{GMm}{r^2}\frac{x}{r}
 =\frac{2GMm\,x}{(x^2+L^2)^{3/2}}.
\]
求导：
\[
 F'(x)=2GMm\frac{L^2-2x^2}{(x^2+L^2)^{5/2}}.
\]
令 $F'(x)=0$，得 $x=L/\sqrt2$。由于 $F(0)=0$ 且 $x\to\infty$ 时 $F\to0$，这必为
最大值。因此在垂直平分线两侧各有一个对称位置：
\[
 \boxed{|x|=\frac{L}{\sqrt2}},
 \qquad
 \boxed{F_{\max}=\frac{4GMm}{3\sqrt3\,L^2}}.
\]
\end{xesolution}

\Needspace{0.43\textheight}
\begin{origexamplebox}{例 1.9}
\begin{enumerate}[label=(\alph*),leftmargin=2.2em]
 \item 从坐标原点以给定初速度 $v_0$ 沿不同角度抛出物体。证明能击中的目标与因初速度
 不足而无法击中的目标之间，分界线满足
 \[
  y=\frac{v_0^2}{2g}-\frac{g}{2v_0^2}x^2.
 \]
 \item 平面上有一个半径为 $R$ 的球体，希望从地面斜向上抛出物体并击中球体顶部。
 在命中球顶以前不接触球体的条件下，求抛出速度的下界，并判断最小值是否能取到；
 再讨论允许提前相切但不穿入球体时的最小抛出速度、抛出角及抛出点到球底部的距离。
\end{enumerate}
\end{origexamplebox}

\begin{xesolution}
\textbf{(a)} 设 $q=\tan\theta$，抛物线轨迹为
\[
 y=xq-\frac{gx^2}{2v_0^2}(1+q^2).
\]
对于给定横坐标 $x$，令右端对 $q$ 取最大值：
\[
 \frac{\partial y}{\partial q}=x-\frac{gx^2}{v_0^2}q=0,\qquad
 q=\frac{v_0^2}{gx}.
\]
代回即得所有同速抛物线的包络线
\[
 \boxed{y_{\max}(x)=\frac{v_0^2}{2g}-\frac{g}{2v_0^2}x^2}.
\]
包络线下方的目标至少对应一个发射角，线上只有一个临界角，线上方不能击中。

\textbf{(b)} 取球底部为原点，球面截线满足
\[
 x^2+(y-R)^2=R^2,
\]
目标是球顶 $(0,2R)$。令抛出点在 $(-d,0)$，以
\[
 X=\frac{x}{R},\quad Y=\frac{y}{R},\quad D=\frac dR,
 \quad Q=\tan\theta,
 \quad A=\frac{gR}{2v_0^2}(1+Q^2)
\]
作无量纲化，则轨迹可写为
\[
 Y=Q(X+D)-A(X+D)^2.
\]
它通过球顶 $(0,2)$，故
\begin{equation}\label{eq:sphere-top}
 QD-AD^2=2.
\end{equation}
先研究允许提前相切的边界轨迹，再检查它是否符合严格不接触的条件。设接触点
横坐标为 $s$，并记 $r=\sqrt{1-s^2}$，则上半圆在该点的纵坐标和斜率分别为
$1+r$ 与 $-s/r$。相切条件为
\begin{equation}\label{eq:sphere-tangent-point}
 Q(s+D)-A(s+D)^2=1+r,
\end{equation}
\begin{equation}\label{eq:sphere-tangent-slope}
 Q-2A(s+D)=-\frac{s}{r}.
\end{equation}
把临界抛物线改写成以相切点为展开中心的形式。因为它在
$X=s$ 处的函数值为 $1+r$、斜率为 $-s/r$，且二次项系数为 $-A$，所以
\[
 Y(X)=1+r-\frac{s}{r}(X-s)-A(X-s)^2.
\]
令 $X=0$，利用轨迹通过球顶 $Y=2$，得到
\[
 2=1+r+\frac{s^2}{r}-As^2.
\]
因此
\[
 A=\frac{r-1+s^2/r}{s^2}
 =\frac{1-r}{s^2r}
 =\frac{1}{r(1+r)},
\]
最后一步使用了 $s^2=1-r^2$。

接下来求发射方向。记
\[
 u=s+D>0,
\]
它是抛出点 $X=-D$ 到相切点 $X=s$ 的水平距离。把 $X=-D$、$Y=0$ 代入以相切点为中心的抛物线，得到
\[
 0=1+r+\frac{s}{r}u-Au^2.
\]
代入 $A=1/[r(1+r)]$ 并整理：
\[
 u^2-s(1+r)u-r(1+r)^2=0.
\]
取 $u>0$ 的根，
\[
 u=\frac{1+r}{2}\left(s+\sqrt{s^2+4r}\right).
\]
在发射点，轨迹斜率就是 $Q=\tan\theta$。由上面的相切点形式求导并令 $X=-D$，
\[
 Q=-\frac{s}{r}+2Au
 =\frac{\sqrt{s^2+4r}}{r}
 =\frac{\sqrt{1+4r-r^2}}{r}.
\]
这同时说明所取的是向右上方发射的 $Q>0$ 分支。于是
\[
 1+Q^2=\frac{1+4r}{r^2}.
\]
再代入 $A=gR(1+Q^2)/(2v_0^2)$，临界速度便成为只含 $r$ 的函数：
\[
 \frac{v_0^2}{gR}=\frac{1+Q^2}{2A}
 =\frac{(r+1)(4r+1)}{2r},\qquad 0<r\le1.
\]
对 $r$ 求导：
\[
 \frac{\dd}{\dd r}\frac{(r+1)(4r+1)}{2r}
 =\frac{(2r-1)(2r+1)}{2r^2}.
\]
最小值在 $r=1/2$ 处取得。此时
\[
 s=-\frac{\sqrt3}{2},\qquad A=\frac43,\qquad
 Q=\sqrt{11},\qquad
 D=\frac{\sqrt3(1+\sqrt{33})}{8}.
\]
因此
\[
 \boxed{v_*=3\sqrt{\frac{gR}{2}}},\qquad
 \boxed{\theta_0=\arctan\sqrt{11}},
\]
\[
 \boxed{d_0=\frac{\sqrt3(1+\sqrt{33})}{8}R}.
\]
\soldiagOneNine
在这组参数下，抛物线在球体左上方
$(x/R,y/R)=(-\sqrt3/2,3/2)$ 处与球面相切，随后到达球顶；因此解图中的“相切”发生在
球顶之前，而不是球顶本身。

还需证明没有另一类轨迹给出更小速度。将任意从左侧到达球顶的轨迹写成
$Y=2-qX-AX^2$。为避免在球顶左邻域穿入球面，必须有 $q\ge0$；
在球顶的水平、竖直速度分别为 $w_x>0$、$-qw_x$，而 $A=gR/(2w_x^2)$。
由竖直运动的速度关系，初速度平方为
\[
 \frac{v_0^2}{gR}=4+\frac{1+q^2}{2A}.
\]
令 $z=-X\in(0,1]$。要求轨迹不低于球的上半圆，等价于
\[
 q\ge h_A(z):=Az-\frac{z}{1+\sqrt{1-z^2}}.
\]
若 $A\le1/2$，可取 $q=0$，但速度平方不小于 $5gR$。
若 $A>1/2$，令 $r=\sqrt{1-z^2}$，直接求导得
\[
 h_A'(z)=A-\frac1{r(1+r)}.
\]
右端随 $z$ 严格减小，由正变负，因此 $h_A$ 恰有一个最大值，
最小可行 $q$ 正是在该处取等号，即前面已经求出的相切轨迹。
这就补足了对所有可行轨迹的比较，而不是只在相切轨迹内部求极值。

\textbf{严格不接触时只有下确界。}等号轨迹提前碰到球面，不能作为严格条件下的解。
固定 $A=4/3$，取 $q=1/\sqrt3+\varepsilon$，其中 $\varepsilon>0$。
在每个 $X<0$ 处，轨迹都比临界轨迹高出 $-\varepsilon X$，仍通过球顶；
令 $\varepsilon\to0^+$，初速度从上方趋于 $v_*$。所以
\[
 \boxed{\inf v_0=v_*,\quad\text{严格不提前接触时最小值不能取到。}}
\]
若允许提前相切而不穿入，才有 $v_{\min}=v_*$，相应角度和距离为上式的 $\theta_0,d_0$。
$\sqrt{gR}$ 具有速度量纲，$d_0$ 与 $R$ 成正比，符合几何相似性。
\end{xesolution}

\subsection{平衡位置及其稳定性}

对一维保守系统，力与势能满足 $F=-\dd U/\dd x$，所以平衡位置首先要求 $U'(x)=0$。
局部极小点受到小扰动后会出现指向平衡点的回复力，因而是稳定平衡；局部极大点则是不稳定平衡。
必须注意：没有阻尼时，物体从局部极小点附近释放后通常会在其两侧振动，并不会自行停在最低点；
只有存在摩擦等耗散时，机械能逐渐减少，系统才可能最终趋近较低势能的稳定平衡。过去常由受力分析
判断平衡；利用势能曲线也能得到相同答案，而且系统越复杂，能量方法往往越简便。

\potentialDiagram

右侧势能曲线中，$S_1,S_2$ 是局部极小点，对应稳定平衡；$U$ 是局部极大点，对应不稳定平衡。
图线其余斜率不为零的位置受力不为零，不能作为静止平衡位置。

以单摆为例，长度为 $l$ 的轻杆一端与 $O$ 点光滑连接，另一端连接质量为 $m$ 的质点。
若以摆竖直向下为 $\theta=0$，系统的重力势能为
\begin{equation}
 E_p=-mgl\cos\theta.
\end{equation}
势能取稳定值的条件为
\begin{equation}
 \frac{\dd E_p}{\dd\theta}=mgl\sin\theta=0,
\end{equation}
故 $\theta=0,\pi$ 都是平衡位置。再求二阶导数
\begin{equation}
 \frac{\dd^2E_p}{\dd\theta^2}=mgl\cos\theta.
\end{equation}
在 $\theta=0$ 处二阶导数为正，是稳定平衡；在 $\theta=\pi$ 处二阶导数为负，是不稳定
平衡。更一般地，势能曲线的局部极小点是稳定平衡，局部极大点是不稳定平衡；导数不为
零的点不能平衡，二阶导数为零的平衡位置须作进一步判断。

\Needspace{0.37\textheight}
\begin{origexamplebox}{例 1.10}
如图，在固定点 $A$ 上系有一根质量可忽略、原长为 $h$、劲度系数为 $k$ 的弹簧，另一
端连接质量为 $m$ 的重物。重物被限制在竖直光滑直线上，$A$ 点到约束直线的距离恰为
$h$。求平衡时重物到 $B$ 点的距离 $x$。
\probdiagOneTen
\end{origexamplebox}

\begin{xesolution}
当重物下降 $x$ 时，弹簧长度
\[
 \ell=\sqrt{h^2+x^2},
\]
伸长量为 $\ell-h$。取 $B$ 点处的重力势能为零，系统势能为
\[
 U(x)=\frac12k(\sqrt{h^2+x^2}-h)^2-mgx.
\]
求导：
\[
 U'(x)=k(\ell-h)\frac{x}{\ell}-mg.
\]
平衡条件为
\[
 \boxed{kx\left(1-\frac{h}{\sqrt{h^2+x^2}}\right)=mg}.
\]
这已经给出 $x$ 的确定方程。若希望写成无根式的多项式，令
$\lambda=mg/(kh)$、$z=x/h$，则
\[
 z\left(1-\frac1{\sqrt{1+z^2}}\right)=\lambda,
\]
等价于
\[
 \boxed{z^4-2\lambda z^3+\lambda^2z^2-2\lambda z+\lambda^2=0},
\]
并须取满足 $z>\lambda$ 的正根，以排除平方带来的伪根。也可以看出
\[
 U''(x)=k\left(1-\frac{h^3}{(h^2+x^2)^{3/2}}\right)>0\quad(x>0),
\]
所以该平衡位置唯一且稳定。一般参数下没有更简洁的初等显式根式，保留隐式方程
更能体现其物理结构。
\end{xesolution}

\Needspace{0.36\textheight}
\begin{origexamplebox}{例 1.11}
如图，锥角为 $\alpha$ 的光滑圆锥体上放着一段质量为 $m$、等效劲度系数为 $k$、原长为
$L_0$ 的弹簧，弹簧始终保持水平。求平衡时弹簧到锥体顶端的距离 $x$。
\probdiagOneEleven
\end{origexamplebox}

\begin{xesolution}
题图中的 $\alpha$ 是圆锥母线与轴线的夹角。距顶端竖直距离为 $x$ 的水平截面半径为
$r=x\tan\alpha$，所以环形弹簧长度
\[
 \ell=2\pi r=2\pi x\tan\alpha.
\]
取顶点高度处为重力势能零点，系统势能为
\[
 U(x)=-mgx+\frac12k(2\pi x\tan\alpha-L_0)^2.
\]
令 $U'(x)=0$：
\[
 -mg+k(2\pi x\tan\alpha-L_0)(2\pi\tan\alpha)=0.
\]
因此
\[
 \boxed{x=\frac{L_0}{2\pi\tan\alpha}
 +\frac{mg}{4\pi^2k\tan^2\alpha}}.
\]
又因
$U''(x)=k(2\pi\tan\alpha)^2>0$，该平衡位置稳定。这里的 $k$ 是整段环形弹簧在周长方向的等效劲度系数。
\end{xesolution}

\Needspace{0.37\textheight}
\begin{origexamplebox}{例 1.12}
两根质量可忽略、长度为 $l$ 的杆在顶部无摩擦连接，质量为 $m$ 的重物 $A$ 放在连接
处。两杆底部放在光滑水平地面上，底端之间由原长为零、劲度系数为 $k$ 的弹簧连接。
求平衡时两杆与地面的夹角 $\theta$。
\probdiagOneTwelve
\end{origexamplebox}

\begin{xesolution}
由对称性，两杆与地面的夹角相同。重物高度为 $l\sin\theta$，两杆底端距离即弹簧长度为
$2l\cos\theta$。系统势能
\[
 U(\theta)=mgl\sin\theta+\frac12k(2l\cos\theta)^2.
\]
求导：
\[
 U'(\theta)=mgl\cos\theta-4kl^2\sin\theta\cos\theta
 =l\cos\theta(mg-4kl\sin\theta).
\]
非退化平衡位置满足
\[
 \boxed{\sin\theta=\frac{mg}{4kl}}.
\]
作为 $0<\theta<\pi/2$ 的内点解，存在条件是 $0<mg<4kl$。在该点
\[
 U''(\theta)=-4kl^2\cos^2\theta<0,
\]
故这一内点平衡是\textbf{不稳定平衡}：角度稍增或稍减，系统都会继续偏离。两个几何
边界 $\theta=0$ 与 $\theta=90^\circ$ 中哪一个是较低势能状态，取决于 $2kl^2$ 与
$mgl$ 的大小；实际能否到达边界，还须结合杆、铰链和重物的接触约束判断。
特别地，$mg=4kl$ 时该解到达 $\theta=\pi/2$，上面的二阶导数为零，不能继续用
$U''<0$ 判断。令 $\delta=\pi/2-\theta$，展开得
$U=4kl^2-\tfrac12kl^2\delta^4+O(\delta^6)$。
若允许向 $\delta>0$ 偏离，这仍是退化的单侧极大值，而不是稳定内点。
\end{xesolution}

\subsection{做展开}

在光滑函数图像上，曲线在某点附近与该点的切线很接近。设函数 $f(x)$ 在 $x_0$ 处的
取值为 $f(x_0)$，切线斜率为 $f'(x_0)$，切线方程为
\begin{equation}
 y=f(x_0)+f'(x_0)(x-x_0).
\end{equation}
当 $x$ 与 $x_0$ 足够接近时，可以近似写为
\begin{equation}
 f(x)\simeq f(x_0)+f'(x_0)(x-x_0).
\end{equation}
若令 $\Delta f=f(x)-f(x_0)$、$\Delta x=x-x_0$，则
\begin{equation}
 \Delta f\simeq f'(x_0)\Delta x.
\end{equation}
这说明连续函数在某点附近的函数值变化量，近似正比于自变量的变化量，比例系数就是
函数在该点的导数。

下面令 $\varepsilon$ 表示无量纲小量，且 $|\varepsilon|\ll1$。常见的小量展开有
\[
 \sin\varepsilon\simeq\tan\varepsilon\simeq\varepsilon,
\]
\[
 (1+\varepsilon)^n\simeq1+n\varepsilon,\qquad
 \e^{\varepsilon}\simeq1+\varepsilon,\qquad
 \ln(1+\varepsilon)\simeq\varepsilon.
\]
它们在 $\varepsilon\to0$ 时与原函数越来越接近。若实际变化量 $\Delta x$ 具有长度量纲，
应先用问题中的特征长度 $L$ 构造 $\varepsilon=\Delta x/L$，不能直接把有量纲的 $\Delta x$ 与数字 $1$ 比较。
小量展开在研究物理系统某点附近的运动行为时非常有效。

\begin{origexamplebox}{例 1.13}
相对论中运动物体的能量随速度的关系不是经典力学中的动能 $\frac12mv^2$，而是
\[
 E(v)=\frac{mc^2}{\sqrt{1-v^2/c^2}}.
\]
证明当物体速度远小于光速时，相对论性能量与经典力学动能给出的结果一致。
\end{origexamplebox}

\begin{xesolution}
令 $\varepsilon=v^2/c^2\ll1$，使用二项式展开
\[
 (1-\varepsilon)^{-1/2}
 =1+\frac12\varepsilon+\frac38\varepsilon^2+\cdots.
\]
因此
\[
 E=mc^2+\frac12mv^2+\frac38m\frac{v^4}{c^2}+\cdots.
\]
第一项 $mc^2$ 是与速度无关的静能，不影响机械运动中的能量差；扣除静能后，相对论动能
为
\[
 E_k=E-mc^2=\frac12mv^2+\mathcal O\!\left(\frac{v^4}{c^2}\right).
\]
所以当 $v\ll c$ 时
\[
 \boxed{E_k\simeq\frac12mv^2},
\]
与经典力学一致。严格地说，题设所给 $E(v)$ 是总能量，不是单独的动能；这正是为何须
扣除 $mc^2$。
\end{xesolution}

\begin{origexamplebox}{例 1.14}
相对论中，质量为 $m$ 的质点在恒力 $F$ 作用下由静止出发做直线运动，位移和时间的关系为
\[
 x(t)=\frac{mc^2}{F}
 \left(\sqrt{1+\left(\frac{Ft}{mc}\right)^2}-1\right).
\]
证明在 $\dfrac{Ft}{mc}\ll1$ 时，相对论公式与经典力学一致。
\end{origexamplebox}

\begin{xesolution}
令 $\varepsilon=Ft/(mc)\ll1$。由
$\sqrt{1+\varepsilon^2}=1+\varepsilon^2/2+\order{\varepsilon^4}$，有
\begin{align*}
 x(t)
 &\simeq\frac{mc^2}{F}\cdot\frac12
 \left(\frac{Ft}{mc}\right)^2\\
 &=\frac12\frac{F}{m}t^2.
\end{align*}
经典力学中恒力产生的加速度为 $a=F/m$，故
\[
 \boxed{x(t)\simeq\frac12at^2}.
\]
下一阶修正为 $-F^3t^4/(8m^3c^2)$，说明相对论位移略小于经典结果。
\end{xesolution}

当一个物理量随另一个物理量改变时，在自变量产生微小变化的情况下，函数值的微小变化
也由给定点处的导数决定。对函数 $y(x)$，当 $x$ 从 $x_0$ 变到 $x_0+\Delta x$ 时，
\begin{equation}
 \Delta y\simeq\frac{\dd y}{\dd x}(x_0)\Delta x.
\end{equation}

\incrementDiagram

上式右边含有两个量：自变量的取值 $x_0$ 与其改变量 $\Delta x$。它们共同给出 $x_0$
附近的函数值变化量，且 $\Delta y$ 正比于该点的导数和 $\Delta x$。

\begin{origexamplebox}{例 1.15}
已知抛体运动的射程
\[
 S=\frac{v^2}{g}\sin2\theta,
\]
其中 $v$ 为初速度，$\theta$ 为抛出角。求速度或角度有微小变化时，射程如何变化。
\end{origexamplebox}

\begin{xesolution}
把 $S$ 看成 $v,\theta$ 的二元函数，其一阶微小变化为全微分
\[
 \Delta S\simeq
 \frac{\partial S}{\partial v}\Delta v
 +\frac{\partial S}{\partial\theta}\Delta\theta.
\]
分别求偏导：
\[
 \frac{\partial S}{\partial v}=\frac{2v}{g}\sin2\theta,\qquad
 \frac{\partial S}{\partial\theta}=\frac{2v^2}{g}\cos2\theta.
\]
所以当两者都可能微变时
\[
 \boxed{\Delta S\simeq
 \frac{2v}{g}\sin2\theta\,\Delta v
 +\frac{2v^2}{g}\cos2\theta\,\Delta\theta}.
\]
若只改变速度，取 $\Delta\theta=0$；若只改变角度，取 $\Delta v=0$，其中角度增量须以
弧度计。特别地，
$\theta=45^\circ$ 时角度的一阶影响为零，这正对应最大射程；此时角度误差的主要影响是
二阶量，而速度误差仍产生一阶影响。
\end{xesolution}

\Needspace{0.40\textheight}
\begin{origexamplebox}{例 1.16}
如图，在 $A$ 点有一个粒子发射源，它在上方 $180^\circ$ 的范围内独立、均匀随机地发射出 $N$ 个
粒子。粒子被与 $A$ 相距 $h$ 的壁接收。建立图示坐标系，求在 $x$ 到 $x+\Delta x$
范围内接收到的粒子数的期望值 $\mathbb E[\Delta n]$，其中 $\Delta x>0$。
假设该小区间足够窄，使方向角在其中可作一阶近似；
例如在 $x=0$ 附近要求 $|\Delta x|\ll h$。
\probdiagOneSixteen
\end{origexamplebox}

\begin{xesolution}
“随机”本身并不意味着均匀；本题另行给出的均匀条件表示等角宽的区间具有相同概率。
长度为 $\pi$ 的角区间中，角宽 $\Delta\theta$ 对应概率 $\Delta\theta/\pi$，
故接收粒子数的期望为
\[
 \mathbb E[\Delta n]=\frac{N}{\pi}\Delta\theta.
\]
由解图的直角三角形
\[
 \tan\theta=\frac{x}{h},\qquad
 \theta=\arctan\frac{x}{h}.
\]
两条边界射线对应的精确角宽，是两个端点方向角之差，因此
\[
 \mathbb E[\Delta n]=\frac N\pi
 \left[\arctan\frac{x+\Delta x}{h}-\arctan\frac{x}{h}\right].
\]
因此
\[
 \frac{\dd\theta}{\dd x}
 =\frac{h}{h^2+x^2},\qquad
 \Delta\theta\simeq\frac{h}{h^2+x^2}\Delta x.
\]
这里的小区间条件不是 $\Delta x\ll x$，而是角度变化率在区间内近似不变。一个简明的充分条件是
$|\Delta x|\ll\sqrt{h^2+x^2}$；在 $x=0$ 处它退化为 $|\Delta x|\ll h$，因此公式仍然完全适用。
于是
\[
 \boxed{\mathbb E[\Delta n]\simeq
 \frac{Nh}{\pi(h^2+x^2)}\Delta x}.
\]
\soldiagOneSixteen
该结果显示墙上粒子分布不均匀：离 $O$ 点越远，单位长度接收的粒子越少。若把墙的两侧
全部计入，积分平均粒子数密度 $\rho(x)=Nh/[\pi(h^2+x^2)]$ 从 $-\infty$ 到 $+\infty$ 恰好得到
总数 $N$，也为系数提供了检验。
一次实验的实际粒子数为整数，会围绕期望值波动；期望值本身不必是整数。
\end{xesolution}
